\section{Experimental Setup}\label{s:experimental-setup}
The evaluation is organized around four research questions.

\parab{RQ1: Conversation quality.}
How does the presence of profile-based resistant clients influence overall conversation quality? Additionally, how does performance vary across different resistance categories?

\parab{RQ2: Diagnostic quality.}
How does the presence of profile-based resistant clients influence overall diagnostic quality? Additionally, how does performance vary across different resistance categories?

\parab{RQ3: Relationship between conversation and diagnosis.}
What is the association between conversation quality and diagnostic quality?

\parab{RQ4: Conversation characteristics.}
How does the presence of profile-based resistant clients influence conversation characteristics? Additionally, how do conversation characteristics vary across different resistance categories?

RQ1 evaluates the quality and clinical usefulness of the generated conversations, while RQ2 evaluates the resulting diagnostic outputs. RQ3 examines the association between these two aspects. RQ4 describes how resistant-client conditions change observable properties of the patient responses and whether the intended resistance is realized in the generated dialogue.


\subsection{Materials}

\paragraph{Assessment Questionnaire.}
All conversations use the adult DSM-5 Self-Rated Level 1 Cross-Cutting Symptom Measure~\cite{apa2013level1}. The measure contains 23 items covering 13 psychiatric domains, including depression, anger, mania, anxiety, psychosis, sleep problems, cognition, and substance use. In this study, the items provide a fixed assessment structure for the assistant--patient conversation. The same questionnaire is used in every experimental condition so that the assessment content remains constant across cooperative and resistant clients.

\paragraph{Clinical Profiles and Diagnostic Materials.}
The experiment uses five clinical profiles included in the DSM5AgentFlow codebase: depression, anxiety, post-traumatic stress disorder, bipolar disorder, and schizophrenia. Each profile defines the expected disorder together with symptoms and background information for the simulated patient. These profiles provide the reference labels used to evaluate the final diagnoses. The diagnostic agent also has access to the local DSM-5 document collection used by DSM5AgentFlow, including APA disorder fact sheets, when producing its diagnostic report.

\paragraph{Resistance Conditions.}
The resistant-client conditions cover the 12 resistance subcategories summarized in Appendix~\ref{s:appendix-taxonomy}. Each resistance subcategory is tested under three intensity conditions: low, medium, and high. The prompt construction, turn-level activation mechanism, expression guidance, cooldown, and truth-map restriction used to generate these conditions are described in Section~\ref{s:methodology}.


\subsection{Conversation Datasets}

\paragraph{Main Dataset.}
The main experiment follows a controlled factorial design. For each of the five clinical profiles, it includes one cooperative baseline and resistant-client conditions covering 12 resistance subcategories at three intensity settings. Every condition is repeated five times. The resulting dataset contains
\[
5 \times 5 \times (12 \times 3 + 1) = 925
\]
conversations. Each conversation contains the complete assistant--patient dialogue, the final clinical observation and diagnostic report, and the condition metadata used in the evaluation.

\begin{table}[t]
    \centering
    \tabcap{Number of conversations in the main experimental dataset.}
    \label{tab:results-dataset-size}
    \begin{tabular}{lrr}
        \toprule
        \thead{Condition} & \thead{Conversations} & \thead{Runs per condition} \\
        \midrule
        Cooperative baseline & 25 & 5 \\
        Low resistance & 300 & 5 \\
        Medium resistance & 300 & 5 \\
        High resistance & 300 & 5 \\
        \midrule
        Total & 925 & -- \\
        \bottomrule
    \end{tabular}
\end{table}

\paragraph{Exploratory Stress-Test Dataset.}
A separate dataset is used to examine the diagnostic pipeline when resistance is applied throughout the questionnaire. For all subcategories except minimization, externalization, and denial, resistance is activated on every questionnaire turn and each activated questionnaire response uses expression level 3. For these three subcategories, resistance is instead activated on every profile-positive questionnaire item identified by the truth map described in Section~\ref{s:methodology}. The dataset contains one conversation for each combination of five clinical profiles and 12 resistance subcategories, giving 60 conversations in total. This stress test is analyzed separately and is not treated as a fourth intensity level in the main factorial experiment.


\subsection{Evaluation Metrics}

The evaluation is organized around three objectives: conversation quality, diagnostic quality, and resistance realization and behavioral characteristics. Table~\ref{tab:metrics-rq-map} summarizes the measures used for each objective.

\begin{table*}[t]
    \centering
    \small
    \tabcap{Evaluation objectives and corresponding measures.}
    \label{tab:metrics-rq-map}
    \begin{tabular}{p{0.20\textwidth} p{0.50\textwidth} p{0.18\textwidth}}
        \toprule
        \thead{Evaluation objective} & \thead{Measures} & \thead{Research questions} \\
        \midrule
        Conversation quality &
        DSM-5 coverage, clinical relevance, consistency, empathy/professionalism, rubric average, question responsiveness, information yield &
        RQ1 and RQ3 \\
        Diagnostic quality &
        Diagnosis classification accuracy, text-match accuracy, diagnostic justification, classification confidence, diagnostic certainty &
        RQ2 and RQ3 \\
        Resistance realization and behavioral characteristics &
        Resistance fidelity, observed activation rate, patient-answer length, activated-answer length, short-answer rate, nonresponse rate &
        RQ4 \\
        \bottomrule
    \end{tabular}
\end{table*}

\subsubsection{Conversation Quality}

Conversation quality is evaluated from both the overall assessment and the patient responses. The rubric-based measures adopted from DSM5AgentFlow score DSM-5 coverage, clinical relevance, logical consistency, diagnostic justification, and empathy/professionalism on five-point scales. Their mean is reported as the rubric average. Together, these measures capture whether the assessment covers clinically relevant material, remains coherent, and provides evidence for the final diagnostic report.

Two additional five-point measures focus on the information provided by the patient. Question responsiveness measures how directly a patient response addresses the current assessment question. Information yield measures how much clinically usable information the response contributes. Long responses do not receive high scores when their additional detail is clinically irrelevant. Both measures are calculated across all patient turns and separately across activated resistance turns. The all-turn scores describe the complete conversation, while the activated-turn scores isolate the responses in which resistance was expressed.

\subsubsection{Diagnostic Quality}

Diagnostic quality is evaluated by comparing the generated diagnosis with the expected diagnosis in the clinical profile. Classification accuracy compares the primary diagnostic class assigned by the evaluator with the expected profile label. Text-match accuracy checks whether the expected disorder name or a recognized variant appears anywhere in the diagnostic report. These measures are reported separately because a report may mention the expected disorder while assigning a different primary diagnosis.

Diagnostic justification, adopted from the DSM5AgentFlow rubric, measures whether the report supports its conclusion with evidence from the conversation. Classification confidence measures how confident the evaluator is when assigning the report to a diagnostic category. Diagnostic certainty measures how confidently the final diagnostic report presents its conclusion, using a five-point scale.

\subsubsection{Resistance Realization and Behavioral Characteristics}

Resistance fidelity is scored on activated turns and measures how closely the generated patient response matches the assigned resistance subcategory. The observed activation rate is calculated directly from the turn-level metadata and shows how often resistance was applied in the generated conversation.

Behavioral characteristics are calculated directly from the patient responses. They include mean patient-answer length, mean activated-answer length, short-answer rate, and nonresponse rate. These measures describe visible differences between resistance categories without requiring an additional LLM judgment. In RQ4, responsiveness and information yield are reported alongside these measures to distinguish response length from clinical usefulness.


\subsection{Implementation Details}
Gemini 2.5 Flash is used for the assistant, patient, and diagnostic agents in all conversation-generation conditions. It is also used for the rubric-based, diagnostic, and client-side LLM evaluations. Using the same model across the cooperative and resistant conditions keeps the model fixed while the client condition changes. The exploratory stress-test conversations are generated with a separate script and are kept separate from the main experimental dataset.

%%% Local Variables:
%%% mode: latex
%%% TeX-master: "../thesis"
%%% End:
